%------------------------------------------------------------------------------%
\documentclass[fleqn,utf8,aspectratio=169,12pt]{beamer}

\usepackage{integracao_e_series}

\title{Volumes por Seções Transversais}

%------------------------------------------------------------------------------%
\begin{document}

\addtitle

%------------------------------------------------------------------------------%
\section{Volume por Seções Transversais}

%------------------------------------------------------------------------------%
\framedgraphic{Volume por Seções Transversais}{volume-secoes-01}{}
\framedgraphic{Volume por Seções Transversais}{volume-secoes-02}{}
\framedgraphic{Volume por Seções Transversais}{volume-secoes-03}{}

%------------------------------------------------------------------------------%
\begin{frame}{Volume por Seções Transversais}
  \emph{$S$}    \tabto{10mm} um sólido

  \pause
  \vspace{\baselineskip}
  \emph{$P_x$}  \tabto{10mm} plano passando por $x$ e perpendicular ao
                             eixo $x$ para $a \leq x \leq b$

  \pause
  \vspace{\baselineskip}
  \emph{$A(x)$} \tabto{10mm} área da secção transversal de $S$ em $P_x$

  \pause
  \vspace{\baselineskip}
  O volume de $S$ é
  \[
    V = \int_a^b A(x)dx
  \]
\end{frame}

%------------------------------------------------------------------------------%
\begin{frame}{Passo a Passo}
\begin{enumerate}[<+->]
  \setlength\itemsep{1em}
  \item Faça o \emph{esboço} do sólido e uma seção transversal típica
  \item Encontre uma \emph{expressão para $A(x)$}
  \item Determine os \emph{limites de integração}
  \item \emph{Integre} $A(x)$ para determinar o volume
\end{enumerate}
\end{frame}

%------------------------------------------------------------------------------%
\section{Exemplos}

%------------------------------------------------------------------------------%
\subsection{Exemplo 1}

%------------------------------------------------------------------------------%
\begin{frame}{Exemplo 1}
  Determine o volume da pirâmide de base quadrada
  com lado \SI{3}{\meter} e \SI{3}{\meter} de altura

  \pause
  \vspace{\baselineskip}
  A seção transversal da pirâmide, perpendicular a altura e a
  $x$ metros abaixo do vértice, é um quadrado com $x$ metros de lado.
\end{frame}

%------------------------------------------------------------------------------%
\begin{frame}{Exemplo 1 -- Volume}
\onslide<+->{\addgraphic{7}{7,2.5}{volume_piramide}}
\begin{align*}
  \onslide<+->{V & = \int_0^3 A(x) dx            \\[2mm]}
  \onslide<+->{  & = \int_0^3 x^2 dx             \\[2mm]}
  \onslide<+->{  & = \frac{x^3}{3} \evalat{0}{3} \\[2mm]}
  \onslide<+->{  & = \frac{27}{3}                \\[2mm]}
  \onslide<+->{  & = 9 \,\si{\meter^3}}
\end{align*}
\end{frame}

%------------------------------------------------------------------------------%
\subsection{Exemplo 2}

%------------------------------------------------------------------------------%
\begin{frame}{Exemplo 2}
  Encontre o volume de uma esfera de raio $r$
\end{frame}

%------------------------------------------------------------------------------%
\begin{frame}{Exemplo 2 -- Área da seção transversal}
  \addgraphic{6.5}{8,2}{volume_esfera}

  \pause
  $P_x$ perpendicular ao eixo $x$

  \pause
  \vspace{\baselineskip}
  Raio do círculo $y=\sqrt{r^2-x^2}$

  \pause
  \vspace{\baselineskip}
  Área da seção transversal
  \[
    A(x)
    = \pi y^2
    \pause
    = \pi (r^2-x^2)
  \]
\end{frame}

%------------------------------------------------------------------------------%
\begin{frame}{Exemplo 2 -- Volume}
  \begin{align*}
    \onslide<+->{V & = \int_{-r}^{r}A(x) \,dx                  &  \\}
    \onslide<+->{  & = \int_{-r}^{r}\pi \left(r^2-x^2\right)dx &  \\}
    \onslide<+->{  & = 2 \pi \int_{0}^{r} r^2-x^2dx            &  & \text{simetria} \hspace{60mm}} ~ \\
    \onslide<+->{  & = 2 \pi \left( r^2 x-\frac{x^3}{3}\right) \evalat{0}{r} & & \text{TFC 2} \\}
    \onslide<+->{  & = 2 \pi \left( r^3-\frac{r^3}{3}\right)                  \\}
    \onslide<+->{  & = \frac{4}{3} \pi r^3                                  }
  \end{align*}
\end{frame}

%------------------------------------------------------------------------------%
\section{Lista Mínima}

%------------------------------------------------------------------------------%
\begin{frame}{Lista Mínima}
  Estudar a Seção 3.3 da Apostila

  Exercícios:

  \vfill
  Atenção: A prova é baseada no livro, não nas apresentações
\end{frame}

%------------------------------------------------------------------------------%
\end{document}
%------------------------------------------------------------------------------%
